% !TeX program = lualatex
% ================================================================
% Polish Tier 3 Vocabulary Version of FactorsMultiplesPrimesQuickQuestions
%
% Generated by the server using the following settings:
%
%   language            Polish (pol)
%   text direction      left to right
%   font                Noto Serif
%   fallback font       Noto Serif
%   built from          latex/quickQuestions/factorsMultiplesPrimesQuickQuestions.tex
%   vocabulary key      latex/quickQuestions/factorsMultiplesPrimesQuickQuestions/L2/pol/factorsMultiplesPrimesQuickQuestions_polishKey.tex
%   tier 3 vocabulary   green   RGB 0 128 0
%   English text        black   RGB 0 0 0
%   Polish text         purple  RGB 112 48 160
%   layout macros       version 3
%
% Please don't edit any information in this comment block - the server
% compares it against the current settings and rebuilds this file when
% they differ, so changes here would be overwritten. However, feel free
% to make updates to the LaTeX below if you want to edit it.
% ================================================================

\documentclass[aspectratio=169]{beamer}
% ================================================================
% Copied from the English sheet this was translated from, so that
% anything it sets up for itself still works here
% ================================================================
\usepackage{amsmath}
\usepackage{amssymb}
\usepackage{tikz}
\usetikzlibrary{shapes.geometric}

% ================================================================
% Quick Questions deck: factors, multiples and primes.
%
% Each topic opens with five true/false statements and then runs ten
% quick questions that get gradually harder. Every question gets two
% slides: the question on its own for the mini whiteboards, then the
% same question again with the solution underneath in red.
%
% Two topics are built differently. "Is it prime?" runs every whole
% number from 1 to 100 on its own slide, then lists that number's
% factors and says why it is or is not prime. "HCF and LCM from lists"
% builds each answer over five slides, adding the factor lists, the
% multiples and then the rings round the HCF and the LCM one slide at
% a time.
% ================================================================

\setbeamertemplate{navigation symbols}{}

% A link in the bottom left of every slide back to the contents, so a
% deck this long can be jumped around in during a lesson.
\setbeamertemplate{footline}{%
  \hbox to\paperwidth{%
    \hspace{1em}%
    \hyperlink{contents}{%
      \color{gray}\scriptsize$\leftarrow$~Back to contents}%
    \hfil}%
  \vskip4pt%
}

% Topic divider.
\newcommand{\qtopic}[1]{%
  \section{#1}%
  \begin{frame}
    \vfill
    \begin{center}
      {\usebeamercolor[fg]{structure}\Huge\bfseries #1}
    \end{center}
    \vfill
  \end{frame}%
}

% \qq[<question size>][<solution size>]{<question>}{<solution>}
% The two optional sizes drop the type for the wordier questions and
% the longer solutions, so nothing runs off the slide.
\NewDocumentCommand{\qq}{O{\Huge}O{\LARGE}+m+m}{%
  \begin{frame}
    \vfill
    \begin{center}
      \begin{minipage}{0.92\textwidth}
        \centering #1 #3
      \end{minipage}
    \end{center}
    \vfill
  \end{frame}%
  \begin{frame}
    \vfill
    \begin{center}
      \begin{minipage}{0.92\textwidth}
        \centering #1 #3
        \par\vspace{0.9em}
        {\color{red}#2 #4}
      \end{minipage}
    \end{center}
    \vfill
  \end{frame}%
}

% \tf[<statement size>][<answer size>]{<statement>}{<answer>}
% As \qq, but with a standing "True or false?" heading above the
% statement. The statements hunt for the usual misconceptions, so the
% answer always says why, not just which.
\NewDocumentCommand{\tf}{O{\LARGE}O{\Large}+m+m}{%
  \begin{frame}
    \vfill
    \begin{center}
      \begin{minipage}{0.92\textwidth}
        \centering
        {\usebeamercolor[fg]{structure}\Large\bfseries True or false?}
        \par\vspace{1em}
        #1 #3
      \end{minipage}
    \end{center}
    \vfill
  \end{frame}%
  \begin{frame}
    \vfill
    \begin{center}
      \begin{minipage}{0.92\textwidth}
        \centering
        {\usebeamercolor[fg]{structure}\Large\bfseries True or false?}
        \par\vspace{1em}
        #1 #3
        \par\vspace{0.9em}
        {\color{red}#2 #4}
      \end{minipage}
    \end{center}
    \vfill
  \end{frame}%
}

% \prq{<number>}{<factors>}{<why it is or is not prime>}
% The "Is it prime?" drill: the number alone in a large centred font,
% then the same number with all of its factors and the reason in red.
% The number is scaled rather than set with \fontsize, because the
% sans font stops at 51.59pt and larger sizes are silently substituted.
\NewDocumentCommand{\prq}{m m +m}{%
  \begin{frame}
    \vfill
    \begin{center}
      \scalebox{3.6}{\Huge\bfseries #1}
    \end{center}
    \vfill
  \end{frame}%
  \begin{frame}
    \vfill
    \begin{center}
      \scalebox{1.9}{\Huge\bfseries #1}
      \par\vspace{0.8em}
      \begin{minipage}{0.9\textwidth}
        \centering\color{red}
        {\large Factors of $#1$: #2}
        \par\vspace{0.7em}
        {\large #3}
      \end{minipage}
    \end{center}
    \vfill
  \end{frame}%
}

% Rings drawn round a number in a list. The ring is set in a box the
% width of the number itself, so the list is spaced as though it were
% ordinary text and nothing moves on the slide where the ring appears;
% the ring itself spills over the gaps either side, as a ring drawn by
% hand would.
\tikzset{
  ringed/.style={ellipse, inner sep=1pt, draw=red, line width=0.9pt},
  unringed/.style={ellipse, inner sep=1pt, draw=none},
}
\newlength{\ringwd}
\newcommand{\ringnode}[2]{%
  \settowidth{\ringwd}{$#2$}%
  \addtolength{\ringwd}{3pt}% a little clearance for the ring to sit in
  \makebox[\ringwd][c]{%
    \tikz[baseline=(n.base)]{\node[#1] (n) {$#2$};}%
  }%
}
\newcommand{\hcfmark}[1]{\alt<4->{\ringnode{ringed}{#1}}{\ringnode{unringed}{#1}}}
\newcommand{\lcmmark}[1]{\alt<5->{\ringnode{ringed}{#1}}{\ringnode{unringed}{#1}}}

% \hcflcm{<a>}{<b>}{<factors of a>}{<factors of b>}
%        {<multiples of a>}{<multiples of b>}
% Five slides: the question, the two factor lists, the two lists of
% multiples, the ring round the HCF, then the ring round the LCM.
% Everything added after the question slide is in red.
% The frame is top aligned so that the question stays put as the lists
% are added underneath it.
\NewDocumentCommand{\hcflcm}{m m +m +m +m +m}{%
  \begin{frame}<1-5>[t]
    \vspace{0.8em}
    \begin{center}
      {\LARGE Find the HCF and LCM of $#1$ and $#2$.}
    \end{center}
    \vspace{1.6em}
    \begin{columns}[T]
      \begin{column}{0.40\textwidth}
        \only<2->{\color{red}\large
          Factors of $#1$:\\[0.35em]
          #3
          \par\vspace{1em}
          Factors of $#2$:\\[0.35em]
          #4}
      \end{column}
      \begin{column}{0.58\textwidth}
        \only<3->{\color{red}\large
          Multiples of $#1$:\\[0.35em]
          #5
          \par\vspace{1em}
          Multiples of $#2$:\\[0.35em]
          #6}
      \end{column}
    \end{columns}
    \vfill
  \end{frame}%
}

% Contents-slide bullets. Green then orange sets the difficulty band;
% blue marks the two topics that work through a method a slide at a
% time rather than asking for a quick answer.
\newcommand{\tocbullet}[1]{\textcolor{#1}{$\bullet$}}
\setbeamertemplate{section in toc}{%
  \ifcase\inserttocsectionnumber
    \or\tocbullet{green!55!black}% 1 Factors
    \or\tocbullet{green!55!black}% 2 Multiples
    \or\tocbullet{green!55!black}% 3 Primes
    \or\tocbullet{blue}%          4 Is it prime?
    \or\tocbullet{orange}%        5 Prime factor decomposition
    \or\tocbullet{blue}%          6 HCF and LCM from lists
    \or\tocbullet{orange}%        7 HCF and LCM by prime factor decomposition
    \or\tocbullet{orange}%        8 Square numbers from prime factors
    \or\tocbullet{orange}%        9 Problem solving with HCF and LCM
  \fi
  ~\inserttocsection%
}

\title{Factors, Multiples and Primes: Quick Questions}
\subtitle{Mini whiteboard questions, with solutions}
\author{}
\date{}

% Characters this language's font has not got are borrowed from another,
% rather than being silently dropped - see docs/translated-sheet-latex.md
\directlua{%
  if luaotfload and luaotfload.add_fallback then
    luaotfload.add_fallback("eallatinfallback",
      { "Noto Serif:mode=harf;" })
  else
    tex.error("this luaotfload is too old to register a fallback font")
  end
}
\usepackage[english]{babel}
\babelprovide[import]{polish}
\babelfont[polish]{rm}[Renderer=Harfbuzz, BoldFont={Noto Serif}, ItalicFont={Noto Serif}, BoldItalicFont={Noto Serif}, RawFeature={fallback=eallatinfallback}]{Noto Serif}
\babelfont[polish]{sf}[Renderer=Harfbuzz, BoldFont={Noto Serif}, ItalicFont={Noto Serif}, BoldItalicFont={Noto Serif}, RawFeature={fallback=eallatinfallback}]{Noto Serif}
\newcommand{\ealtext}[1]{\foreignlanguage{polish}{#1}}
\newcommand{\ealtextblock}[1]{{\raggedright\foreignlanguage{polish}{#1}\par}}
% ================================================================
% EAL layout helpers
% ================================================================
\usepackage{xcolor}

\definecolor{ealkeycolour}{RGB}{0,128,0}
\definecolor{ealenglish}{RGB}{0,0,0}
\definecolor{ealtwocolour}{RGB}{112,48,160}

\newsavebox{\ealboxen}
\newsavebox{\ealboxtr}
\newlength{\ealglosswd}
\newlength{\ealglossraise}
\setlength{\ealglossraise}{1.15em}

% a tier 3 word, in English and in translation alike
\newcommand{\ealkey}[1]{\textcolor{ealkeycolour}{#1}}
\newcommand{\ealkeytr}[1]{\textcolor{ealkeycolour}{\ealtext{#1}}}

% One translated word above one English word. Built from kernel
% primitives, never a tabular, which would break wherever the sheet
% loads array, booktabs or tabularx. The box is as wide as the wider of
% the two, so a long translation pushes its neighbours aside instead of
% overlapping them, and it claims its full height so a table row or a
% TikZ node makes room for it.
\newcommand{\ealgloss}[2]{%
  \sbox{\ealboxen}{\ealkey{#1}}%
  \sbox{\ealboxtr}{\small\textcolor{ealtwocolour}{\ealtext{#2}}}%
  \setlength{\ealglosswd}{\wd\ealboxen}%
  \ifdim\wd\ealboxtr>\ealglosswd\setlength{\ealglosswd}{\wd\ealboxtr}\fi
  \leavevmode
  \makebox[\ealglosswd][c]{%
    \makebox[0pt][c]{\usebox{\ealboxen}}%
    \makebox[0pt][c]{\raisebox{\ealglossraise}{\usebox{\ealboxtr}}}%
  }%
}

% opens up line spacing around a block of glosses, so the raised
% translations sit clear of the line above
\newenvironment{ealglossed}{\par\linespread{2.1}\selectfont\ignorespaces}{\par}

% a whole translated sentence above its English counterpart. block level,
% so it does not belong inside a tabular cell
\newcommand{\ealpara}[2]{%
  \par\addvspace{0.5\baselineskip}%
  {\color{ealtwocolour}\ealtextblock{#1}}%
  \penalty200
  {\color{ealenglish}#2\par}%
  \addvspace{0.5\baselineskip}%
}

\begin{document}

% ---- EAL helper setup (original custom macros, modified) ----
\setbeamertemplate{navigation symbols}{}

\setbeamertemplate{footline}{%
  \hbox to\paperwidth{%
    \hspace{1em}%
    \hyperlink{contents}{%
      \color{gray}\scriptsize$\leftarrow$~Back to contents}%
    \hfil}%
  \vskip4pt%
}

\newcommand{\qtopic}[1]{%
  \begin{frame}
    \vfill
    \begin{center}
      \begin{ealglossed}
      {\usebeamercolor[fg]{structure}\Huge\bfseries #1}
      \end{ealglossed}
    \end{center}
    \vfill
  \end{frame}%
}

\NewDocumentCommand{\qq}{O{\Huge}O{\LARGE}+m+m}{%
  \begin{frame}
    \vfill
    \begin{center}
      \begin{minipage}{0.92\textwidth}
        \centering #1 \begin{ealglossed}#3\end{ealglossed}
      \end{minipage}
    \end{center}
    \vfill
  \end{frame}%
  \begin{frame}
    \vfill
    \begin{center}
      \begin{minipage}{0.92\textwidth}
        \centering #1 \begin{ealglossed}#3\end{ealglossed}
        \par\vspace{0.9em}
        {\color{red}#2 \begin{ealglossed}#4\end{ealglossed}}
      \end{minipage}
    \end{center}
    \vfill
  \end{frame}%
}

\NewDocumentCommand{\tf}{O{\LARGE}O{\Large}+m+m}{%
  \begin{frame}
    \vfill
    \begin{center}
      \begin{minipage}{0.92\textwidth}
        \centering
        {\usebeamercolor[fg]{structure}\Large\bfseries True or false?}
        \par\vspace{1em}
        #1 \begin{ealglossed}#3\end{ealglossed}
      \end{minipage}
    \end{center}
    \vfill
  \end{frame}%
  \begin{frame}
    \vfill
    \begin{center}
      \begin{minipage}{0.92\textwidth}
        \centering
        {\usebeamercolor[fg]{structure}\Large\bfseries True or false?}
        \par\vspace{1em}
        #1 \begin{ealglossed}#3\end{ealglossed}
        \par\vspace{0.9em}
        {\color{red}#2 \begin{ealglossed}#4\end{ealglossed}}
      \end{minipage}
    \end{center}
    \vfill
  \end{frame}%
}

\NewDocumentCommand{\prq}{m m +m}{%
  \begin{frame}
    \vfill
    \begin{center}
      \scalebox{3.6}{\Huge\bfseries #1}
    \end{center}
    \vfill
  \end{frame}%
  \begin{frame}
    \vfill
    \begin{center}
      \scalebox{1.9}{\Huge\bfseries #1}
      \par\vspace{0.8em}
      \begin{minipage}{0.9\textwidth}
        \centering\color{red}
        \begin{ealglossed}
        {\large \ealgloss{Factors}{Dzielniki} of $#1$: #2}
        \par\vspace{0.7em}
        {\large #3}
        \end{ealglossed}
      \end{minipage}
    \end{center}
    \vfill
  \end{frame}%
}

\tikzset{
  ringed/.style={ellipse, inner sep=1pt, draw=red, line width=0.9pt},
  unringed/.style={ellipse, inner sep=1pt, draw=none},
}
\newlength{\ringwd}
\newcommand{\ringnode}[2]{%
  \settowidth{\ringwd}{$#2$}%
  \addtolength{\ringwd}{3pt}%
  \makebox[\ringwd][c]{%
    \tikz[baseline=(n.base)]{\node[#1] (n) {$#2$};}%
  }%
}
\newcommand{\hcfmark}[1]{\alt<4->{\ringnode{ringed}{#1}}{\ringnode{unringed}{#1}}}
\newcommand{\lcmmark}[1]{\alt<5->{\ringnode{ringed}{#1}}{\ringnode{unringed}{#1}}}

\NewDocumentCommand{\hcflcm}{m m +m +m +m +m}{%
  \begin{frame}<1-5>[t]
    \begin{ealglossed}
    \vspace{0.8em}
    \begin{center}
      {\LARGE Find the \ealgloss{HCF}{największy wspólny dzielnik} and \ealgloss{LCM}{najmniejsza wspólna wielokrotność} of $#1$ and $#2$.}
    \end{center}
    \vspace{1.6em}
    \begin{columns}[T]
      \begin{column}{0.40\textwidth}
        \only<2->{\color{red}\large
          \ealgloss{Factors}{Dzielniki} of $#1$:\\[0.35em]
          #3
          \par\vspace{1em}
          \ealgloss{Factors}{Dzielniki} of $#2$:\\[0.35em]
          #4}
      \end{column}
      \begin{column}{0.58\textwidth}
        \only<3->{\color{red}\large
          \ealgloss{Multiples}{Wielokrotności} of $#1$:\\[0.35em]
          #5
          \par\vspace{1em}
          \ealgloss{Multiples}{Wielokrotności} of $#2$:\\[0.35em]
          #6}
      \end{column}
    \end{columns}
    \vfill
    \end{ealglossed}
  \end{frame}%
}

% ---- Title ----
\title{\ealgloss{Factors}{Dzielniki}, \ealgloss{Multiples}{Wielokrotności} and \ealgloss{Primes}{Liczby pierwsze}: Quick Questions}
\subtitle{Mini whiteboard questions, with solutions}
\author{}
\date{}

\begin{frame}
  \begin{ealglossed}
  \titlepage
  \end{ealglossed}
\end{frame}

\begin{frame}[label=contents]{Contents}
  \begin{ealglossed}
  \small
  \begin{itemize}
    \item \ealgloss{Factors}{Dzielniki}
    \item \ealgloss{Multiples}{Wielokrotności}
    \item \ealgloss{Primes}{Liczby pierwsze}
    \item Is it \ealgloss{prime}{liczba pierwsza}?
    \item \ealgloss{Prime factor decomposition}{Rozkład na czynniki pierwsze}
    \item \ealgloss{HCF}{największy wspólny dzielnik} and \ealgloss{LCM}{najmniejsza wspólna wielokrotność} from lists
    \item \ealgloss{HCF}{największy wspólny dzielnik} and \ealgloss{LCM}{najmniejsza wspólna wielokrotność} by \ealgloss{prime factor decomposition}{rozkład na czynniki pierwsze}
    \item \ealgloss{Square numbers}{Liczby kwadratowe} from \ealgloss{prime factors}{czynniki pierwsze}
    \item Problem solving with \ealgloss{HCF}{największy wspólny dzielnik} and \ealgloss{LCM}{najmniejsza wspólna wielokrotność}
  \end{itemize}
  \end{ealglossed}
\end{frame}

%================================================================
\qtopic{\ealgloss{Factors}{Dzielniki}}
%================================================================

\tf{$7$ is a \ealgloss{factor}{dzielnik} of $21$}
  {\textbf{True.} $21=7\times 3$, so $21$ divides exactly by $7$.}

\tf{$1$ is a \ealgloss{factor}{dzielnik} of every whole number}
  {\textbf{True.} Every number $n$ can be written as $1\times n$.}

\tf{$6$ is a \ealgloss{factor}{dzielnik} of $16$}
  {\textbf{False.} The \ealgloss{factors}{dzielniki} of $16$ are $1,2,4,8,16$;\\[0.3em]
   $16\div 6$ is not a whole number.}

\tf{Every number has an even number of \ealgloss{factors}{dzielniki}}
  {\textbf{False.} \ealgloss{Factors}{Dzielniki} usually pair up, but $9$ has just $1,3,9$
   --- three \ealgloss{factors}{dzielniki} --- because $3\times 3$ uses the same \ealgloss{factor}{dzielnik} twice.}

\tf{If $4$ is a \ealgloss{factor}{dzielnik} of a number,\\[0.3em]then $8$ must be a \ealgloss{factor}{dzielnik} of it too}
  {\textbf{False.} $4$ is a \ealgloss{factor}{dzielnik} of $12$, but $8$ is not.}

\qq[\LARGE]{List all the \ealgloss{factors}{dzielniki} of $12$}{$1,\ 2,\ 3,\ 4,\ 6,\ 12$}

\qq[\LARGE]{List all the \ealgloss{factors}{dzielniki} of $20$}{$1,\ 2,\ 4,\ 5,\ 10,\ 20$}

\qq[\LARGE]{List all the \ealgloss{factors}{dzielniki} of $36$}
  {$1,\ 2,\ 3,\ 4,\ 6,\ 9,\ 12,\ 18,\ 36$}

\qq[\LARGE]{How many \ealgloss{factors}{dzielniki} does $16$ have?}
  {$1,\ 2,\ 4,\ 8,\ 16$\\[0.3em]Five \ealgloss{factors}{dzielniki}.}

\qq[\LARGE]{Find the \ealgloss{highest common factor}{największy wspólny dzielnik} of $12$ and $18$}
  {\ealgloss{Factors}{Dzielniki} of $12$: $1,\ 2,\ 3,\ 4,\ \mathbf{6},\ 12$\\[0.3em]
   \ealgloss{Factors}{Dzielniki} of $18$: $1,\ 2,\ 3,\ \mathbf{6},\ 9,\ 18$\\[0.3em]
   $\mathrm{HCF}=6$}

\qq[\LARGE]{Find the \ealgloss{highest common factor}{największy wspólny dzielnik} of $24$ and $36$}
  {\ealgloss{Factors}{Dzielniki} of $24$: $1,\ 2,\ 3,\ 4,\ 6,\ 8,\ \mathbf{12},\ 24$\\[0.3em]
   \ealgloss{Factors}{Dzielniki} of $36$: $1,\ 2,\ 3,\ 4,\ 6,\ 9,\ \mathbf{12},\ 18,\ 36$\\[0.3em]
   $\mathrm{HCF}=12$}

\qq[\LARGE]{List all the \ealgloss{factor pairs}{pary czynników} of $48$}
  {$1\times 48\qquad 2\times 24\qquad 3\times 16$\\[0.4em]
   $4\times 12\qquad 6\times 8$}

\qq[\LARGE]{Which whole number below $30$ has the most \ealgloss{factors}{dzielniki}?}
  {$24$, with eight \ealgloss{factors}{dzielniki}:\\[0.3em]
   $1,\ 2,\ 3,\ 4,\ 6,\ 8,\ 12,\ 24$}

\qq[\LARGE]{$n$ has exactly three \ealgloss{factors}{dzielniki}.\\[0.4em]
Give a possible value of $n$.}
  {$4$ ($1,2,4$), $9$ ($1,3,9$) or $25$ ($1,5,25$).\\[0.3em]
   Any \ealgloss{prime}{liczba pierwsza} squared works.}

\qq[\Large][\Large]{Explain why \ealgloss{square numbers}{liczby kwadratowe} have an odd number of \ealgloss{factors}{dzielniki}}
  {\ealgloss{Factors}{Dzielniki} come in pairs, one either side of the \ealgloss{square root}{pierwiastek}.\\[0.3em]
   In a \ealgloss{square number}{liczba kwadratowa} the \ealgloss{square root}{pierwiastek} pairs with itself,
   so one \ealgloss{factor}{dzielnik} is left unpaired.}

%================================================================
\qtopic{\ealgloss{Multiples}{Wielokrotności}}
%================================================================

\tf{$24$ is a \ealgloss{multiple}{wielokrotność} of $6$}
  {\textbf{True.} $6\times 4=24$.}

\tf{Every \ealgloss{multiple}{wielokrotność} of $10$ is also a \ealgloss{multiple}{wielokrotność} of $5$}
  {\textbf{True.} $10=5\times 2$, so every \ealgloss{multiple}{wielokrotność} of $10$
   already contains a \ealgloss{factor}{dzielnik} of $5$.}

\tf{$35$ is a \ealgloss{multiple}{wielokrotność} of $4$}
  {\textbf{False.} $4\times 8=32$ and $4\times 9=36$, so $35$ is stepped over.}

\tf{Every \ealgloss{multiple}{wielokrotność} of $3$ is also a \ealgloss{multiple}{wielokrotność} of $9$}
  {\textbf{False.} $6$ is a \ealgloss{multiple}{wielokrotność} of $3$ but not of $9$.\\[0.3em]
   The other way round \emph{is} true.}

\tf[\Large]{The \ealgloss{lowest common multiple}{najmniejsza wspólna wielokrotność} of two numbers is always found by
multiplying them together}
  {\textbf{False.} $4\times 6=24$, but the \ealgloss{LCM}{najmniejsza wspólna wielokrotność} of $4$ and $6$ is $12$.\\[0.3em]
   Multiplying only works when the numbers share no \ealgloss{factors}{dzielniki}.}

\qq[\LARGE]{Write down the first five \ealgloss{multiples}{wielokrotności} of $4$}
  {$4,\ 8,\ 12,\ 16,\ 20$}

\qq[\LARGE]{Write down the first five \ealgloss{multiples}{wielokrotności} of $7$}
  {$7,\ 14,\ 21,\ 28,\ 35$}

\qq[\LARGE]{Write down the \ealgloss{multiples}{wielokrotności} of $6$ between $20$ and $50$}
  {$24,\ 30,\ 36,\ 42,\ 48$}

\qq[\LARGE]{Is $126$ a \ealgloss{multiple}{wielokrotność} of $9$?}
  {Yes. $1+2+6=9$, and $9\times 14=126$.}

\qq[\LARGE]{Find the \ealgloss{lowest common multiple}{najmniejsza wspólna wielokrotność} of $4$ and $6$}
  {$4,\ 8,\ \mathbf{12},\ 16,\ \ldots$\\[0.3em]
   $6,\ \mathbf{12},\ 18,\ \ldots$\\[0.3em]
   $\mathrm{LCM}=12$}

\qq[\LARGE]{Find the \ealgloss{lowest common multiple}{najmniejsza wspólna wielokrotność} of $8$ and $12$}
  {$8,\ 16,\ \mathbf{24},\ 32,\ \ldots$\\[0.3em]
   $12,\ \mathbf{24},\ 36,\ \ldots$\\[0.3em]
   $\mathrm{LCM}=24$}

\qq[\LARGE]{Find the \ealgloss{lowest common multiple}{najmniejsza wspólna wielokrotność} of $5$, $6$ and $9$}
  {\ealgloss{LCM}{najmniejsza wspólna wielokrotność} of $5$ and $6$ is $30$.\\[0.3em]
   \ealgloss{LCM}{najmniejsza wspólna wielokrotność} of $30$ and $9$ is $90$.\\[0.3em]
   $\mathrm{LCM}=90$}

\qq[\LARGE]{A number is a \ealgloss{multiple}{wielokrotność} of both $4$ and $10$.\\[0.4em]
What is the smallest it could be?}
  {$20$, the \ealgloss{lowest common multiple}{najmniejsza wspólna wielokrotność} of $4$ and $10$.}

\qq[\Large][\Large]{Two lighthouses flash together now.\\[0.4em]
One flashes every $12$ seconds, the other every $18$ seconds.\\[0.4em]
When do they next flash together?}
  {$\mathrm{LCM}(12,18)=36$\\[0.3em]
   In $36$ seconds.}

\qq[\Large][\Large]{Find the smallest number above $1$ that leaves a remainder
of $1$ when divided by $4$, by $5$ and by $6$}
  {The number must be one more than a common \ealgloss{multiple}{wielokrotność} of $4$, $5$ and $6$.\\[0.3em]
   $\mathrm{LCM}(4,5,6)=60$, so the number is $61$.}

%================================================================
\qtopic{\ealgloss{Primes}{Liczby pierwsze}}
%================================================================

\tf{$1$ is a \ealgloss{prime}{liczba pierwsza} number}
  {\textbf{False.} A \ealgloss{prime}{liczba pierwsza} has exactly two \ealgloss{factors}{dzielniki}, but $1$ has only one.}

\tf{$2$ is the only even \ealgloss{prime}{liczba pierwsza} number}
  {\textbf{True.} Every other even number has $2$ as an extra \ealgloss{factor}{dzielnik},
   so it has more than two \ealgloss{factors}{dzielniki}.}

\tf{$51$ is a \ealgloss{prime}{liczba pierwsza} number}
  {\textbf{False.} $51=3\times 17$.\\[0.3em]
   Its digits add to $6$, so it divides by $3$.}

\tf{There are infinitely many \ealgloss{prime}{liczby pierwsze} numbers}
  {\textbf{True.} The \ealgloss{primes}{liczby pierwsze} never run out,
   although they do thin out as the numbers get larger.}

\tf{If a number ends in $3$, it must be \ealgloss{prime}{liczba pierwsza}}
  {\textbf{False.} $33=3\times 11$ and $63=7\times 9$.}

\qq[\LARGE]{Write down all the \ealgloss{prime}{liczby pierwsze} numbers between $1$ and $10$}
  {$2,\ 3,\ 5,\ 7$}

\qq[\LARGE]{Write down all the \ealgloss{prime}{liczby pierwsze} numbers between $10$ and $20$}
  {$11,\ 13,\ 17,\ 19$}

\qq[\LARGE]{Is $27$ a \ealgloss{prime}{liczba pierwsza} number?\\[0.4em]Explain your answer.}
  {No. $27=3\times 9$, so it has more than two \ealgloss{factors}{dzielniki}.}

\qq[\LARGE]{Write down the first \ealgloss{prime}{liczba pierwsza} number greater than $30$}
  {$31$}

\qq[\LARGE]{Is $91$ a \ealgloss{prime}{liczba pierwsza} number?}
  {No. $91=7\times 13$.}

\qq[\LARGE]{How many \ealgloss{prime}{liczby pierwsze} numbers are there between $20$ and $40$?}
  {$23,\ 29,\ 31,\ 37$\\[0.3em]Four of them.}

\qq[\LARGE]{Is $143$ a \ealgloss{prime}{liczba pierwsza} number?}
  {No. $143=11\times 13$.}

\qq[\LARGE]{Find two \ealgloss{prime}{liczby pierwsze} numbers that add up to $24$}
  {$5+19$ \quad $7+17$ \quad $11+13$}

\qq[\Large][\Large]{To test whether $211$ is \ealgloss{prime}{liczba pierwsza},\\[0.3em]
which numbers do you need to divide by?\\[0.4em]Is it \ealgloss{prime}{liczba pierwsza}?}
  {Only the \ealgloss{primes}{liczby pierwsze} up to $\sqrt{211}\approx 14.5$:\\[0.3em]
   $2,\ 3,\ 5,\ 7,\ 11,\ 13$\\[0.3em]
   None of them divides $211$, so $211$ is \ealgloss{prime}{liczba pierwsza}.}

\qq[\Large][\Large]{Explain why $n^2-1$ can never be \ealgloss{prime}{liczba pierwsza} when $n>2$}
  {$n^2-1=(n-1)(n+1)$\\[0.3em]
   When $n>2$ both brackets are bigger than $1$,
   so the number always has more than two \ealgloss{factors}{dzielniki}.}

%================================================================
\qtopic{Is it \ealgloss{prime}{liczba pierwsza}?}
%================================================================

\prq{1}{$1$}
  {$1$ has only one \ealgloss{factor}{dzielnik}. A \ealgloss{prime}{liczba pierwsza} number has exactly two different \ealgloss{factors}{dzielniki}, $1$ and itself, and for $1$ those would be the same \ealgloss{factor}{dzielnik}, so $1$ is not a \ealgloss{prime}{liczba pierwsza} number.}

\prq{2}{$1$, $2$}
  {$2$ has exactly two \ealgloss{factors}{dzielniki}, $1$ and itself, so $2$ is a \ealgloss{prime}{liczba pierwsza} number. It is the only even \ealgloss{prime}{liczba pierwsza} number, because every other even number has $2$ as a third \ealgloss{factor}{dzielnik}.}

\prq{3}{$1$, $3$}
  {$3$ has exactly two \ealgloss{factors}{dzielniki}, $1$ and itself, so $3$ is a \ealgloss{prime}{liczba pierwsza} number.}

\prq{4}{$1$, $2$, $4$}
  {$4$ has three \ealgloss{factors}{dzielniki}, not two, so $4$ is not a \ealgloss{prime}{liczba pierwsza} number.}

\prq{5}{$1$, $5$}
  {$5$ has exactly two \ealgloss{factors}{dzielniki}, $1$ and itself, so $5$ is a \ealgloss{prime}{liczba pierwsza} number.\\[0.45em]We only need to test $2$ as a \ealgloss{factor}{dzielnik}, because $3\times 3=9$ is bigger than $5$.}

\prq{6}{$1$, $2$, $3$, $6$}
  {$6$ has four \ealgloss{factors}{dzielniki}, not two, so $6$ is not a \ealgloss{prime}{liczba pierwsza} number.}

\prq{7}{$1$, $7$}
  {$7$ has exactly two \ealgloss{factors}{dzielniki}, $1$ and itself, so $7$ is a \ealgloss{prime}{liczba pierwsza} number.\\[0.45em]We only need to test $2$ as a \ealgloss{factor}{dzielnik}, because $3\times 3=9$ is bigger than $7$.}

\prq{8}{$1$, $2$, $4$, $8$}
  {$8$ has four \ealgloss{factors}{dzielniki}, not two, so $8$ is not a \ealgloss{prime}{liczba pierwsza} number.}

\prq{9}{$1$, $3$, $9$}
  {$9$ has three \ealgloss{factors}{dzielniki}, not two, so $9$ is not a \ealgloss{prime}{liczba pierwsza} number.}

\prq{10}{$1$, $2$, $5$, $10$}
  {$10$ has four \ealgloss{factors}{dzielniki}, not two, so $10$ is not a \ealgloss{prime}{liczba pierwsza} number.}

\prq{11}{$1$, $11$}
  {$11$ has exactly two \ealgloss{factors}{dzielniki}, $1$ and itself, so $11$ is a \ealgloss{prime}{liczba pierwsza} number.\\[0.45em]We only need to test $2$ and $3$ as \ealgloss{factors}{dzielniki}, because $5\times 5=25$ is bigger than $11$.}

\prq{12}{$1$, $2$, $3$, $4$, $6$, $12$}
  {$12$ has six \ealgloss{factors}{dzielniki}, not two, so $12$ is not a \ealgloss{prime}{liczba pierwsza} number.}

\prq{13}{$1$, $13$}
  {$13$ has exactly two \ealgloss{factors}{dzielniki}, $1$ and itself, so $13$ is a \ealgloss{prime}{liczba pierwsza} number.\\[0.45em]We only need to test $2$ and $3$ as \ealgloss{factors}{dzielniki}, because $5\times 5=25$ is bigger than $13$.}

\prq{14}{$1$, $2$, $7$, $14$}
  {$14$ has four \ealgloss{factors}{dzielniki}, not two, so $14$ is not a \ealgloss{prime}{liczba pierwsza} number.}

\prq{15}{$1$, $3$, $5$, $15$}
  {$15$ has four \ealgloss{factors}{dzielniki}, not two, so $15$ is not a \ealgloss{prime}{liczba pierwsza} number.}

\prq{16}{$1$, $2$, $4$, $8$, $16$}
  {$16$ has five \ealgloss{factors}{dzielniki}, not two, so $16$ is not a \ealgloss{prime}{liczba pierwsza} number.}

\prq{17}{$1$, $17$}
  {$17$ has exactly two \ealgloss{factors}{dzielniki}, $1$ and itself, so $17$ is a \ealgloss{prime}{liczba pierwsza} number.\\[0.45em]We only need to test $2$ and $3$ as \ealgloss{factors}{dzielniki}, because $5\times 5=25$ is bigger than $17$.}

\prq{18}{$1$, $2$, $3$, $6$, $9$, $18$}
  {$18$ has six \ealgloss{factors}{dzielniki}, not two, so $18$ is not a \ealgloss{prime}{liczba pierwsza} number.}

\prq{19}{$1$, $19$}
  {$19$ has exactly two \ealgloss{factors}{dzielniki}, $1$ and itself, so $19$ is a \ealgloss{prime}{liczba pierwsza} number.\\[0.45em]We only need to test $2$ and $3$ as \ealgloss{factors}{dzielniki}, because $5\times 5=25$ is bigger than $19$.}

\prq{20}{$1$, $2$, $4$, $5$, $10$, $20$}
  {$20$ has six \ealgloss{factors}{dzielniki}, not two, so $20$ is not a \ealgloss{prime}{liczba pierwsza} number.}

\prq{21}{$1$, $3$, $7$, $21$}
  {$21$ has four \ealgloss{factors}{dzielniki}, not two, so $21$ is not a \ealgloss{prime}{liczba pierwsza} number.}

\prq{22}{$1$, $2$, $11$, $22$}
  {$22$ has four \ealgloss{factors}{dzielniki}, not two, so $22$ is not a \ealgloss{prime}{liczba pierwsza} number.}

\prq{23}{$1$, $23$}
  {$23$ has exactly two \ealgloss{factors}{dzielniki}, $1$ and itself, so $23$ is a \ealgloss{prime}{liczba pierwsza} number.\\[0.45em]We only need to test $2$ and $3$ as \ealgloss{factors}{dzielniki}, because $5\times 5=25$ is bigger than $23$.}

\prq{24}{$1$, $2$, $3$, $4$, $6$, $8$, $12$, $24$}
  {$24$ has eight \ealgloss{factors}{dzielniki}, not two, so $24$ is not a \ealgloss{prime}{liczba pierwsza} number.}

\prq{25}{$1$, $5$, $25$}
  {$25$ has three \ealgloss{factors}{dzielniki}, not two, so $25$ is not a \ealgloss{prime}{liczba pierwsza} number.}

\prq{26}{$1$, $2$, $13$, $26$}
  {$26$ has four \ealgloss{factors}{dzielniki}, not two, so $26$ is not a \ealgloss{prime}{liczba pierwsza} number.}

\prq{27}{$1$, $3$, $9$, $27$}
  {$27$ has four \ealgloss{factors}{dzielniki}, not two, so $27$ is not a \ealgloss{prime}{liczba pierwsza} number.}

\prq{28}{$1$, $2$, $4$, $7$, $14$, $28$}
  {$28$ has six \ealgloss{factors}{dzielniki}, not two, so $28$ is not a \ealgloss{prime}{liczba pierwsza} number.}

\prq{29}{$1$, $29$}
  {$29$ has exactly two \ealgloss{factors}{dzielniki}, $1$ and itself, so $29$ is a \ealgloss{prime}{liczba pierwsza} number.\\[0.45em]We only need to test $2$, $3$ and $5$ as \ealgloss{factors}{dzielniki}, because $7\times 7=49$ is bigger than $29$.}

\prq{30}{$1$, $2$, $3$, $5$, $6$, $10$, $15$, $30$}
  {$30$ has eight \ealgloss{factors}{dzielniki}, not two, so $30$ is not a \ealgloss{prime}{liczba pierwsza} number.}

\prq{31}{$1$, $31$}
  {$31$ has exactly two \ealgloss{factors}{dzielniki}, $1$ and itself, so $31$ is a \ealgloss{prime}{liczba pierwsza} number.\\[0.45em]We only need to test $2$, $3$ and $5$ as \ealgloss{factors}{dzielniki}, because $7\times 7=49$ is bigger than $31$.}

\prq{32}{$1$, $2$, $4$, $8$, $16$, $32$}
  {$32$ has six \ealgloss{factors}{dzielniki}, not two, so $32$ is not a \ealgloss{prime}{liczba pierwsza} number.}

\prq{33}{$1$, $3$, $11$, $33$}
  {$33$ has four \ealgloss{factors}{dzielniki}, not two, so $33$ is not a \ealgloss{prime}{liczba pierwsza} number.}

\prq{34}{$1$, $2$, $17$, $34$}
  {$34$ has four \ealgloss{factors}{dzielniki}, not two, so $34$ is not a \ealgloss{prime}{liczba pierwsza} number.}

\prq{35}{$1$, $5$, $7$, $35$}
  {$35$ has four \ealgloss{factors}{dzielniki}, not two, so $35$ is not a \ealgloss{prime}{liczba pierwsza} number.}

\prq{36}{$1$, $2$, $3$, $4$, $6$, $9$, $12$, $18$, $36$}
  {$36$ has nine \ealgloss{factors}{dzielniki}, not two, so $36$ is not a \ealgloss{prime}{liczba pierwsza} number.}

\prq{37}{$1$, $37$}
  {$37$ has exactly two \ealgloss{factors}{dzielniki}, $1$ and itself, so $37$ is a \ealgloss{prime}{liczba pierwsza} number.\\[0.45em]We only need to test $2$, $3$ and $5$ as \ealgloss{factors}{dzielniki}, because $7\times 7=49$ is bigger than $37$.}

\prq{38}{$1$, $2$, $19$, $38$}
  {$38$ has four \ealgloss{factors}{dzielniki}, not two, so $38$ is not a \ealgloss{prime}{liczba pierwsza} number.}

\prq{39}{$1$, $3$, $13$, $39$}
  {$39$ has four \ealgloss{factors}{dzielniki}, not two, so $39$ is not a \ealgloss{prime}{liczba pierwsza} number.}

\prq{40}{$1$, $2$, $4$, $5$, $8$, $10$, $20$, $40$}
  {$40$ has eight \ealgloss{factors}{dzielniki}, not two, so $40$ is not a \ealgloss{prime}{liczba pierwsza} number.}

\prq{41}{$1$, $41$}
  {$41$ has exactly two \ealgloss{factors}{dzielniki}, $1$ and itself, so $41$ is a \ealgloss{prime}{liczba pierwsza} number.\\[0.45em]We only need to test $2$, $3$ and $5$ as \ealgloss{factors}{dzielniki}, because $7\times 7=49$ is bigger than $41$.}

\prq{42}{$1$, $2$, $3$, $6$, $7$, $14$, $21$, $42$}
  {$42$ has eight \ealgloss{factors}{dzielniki}, not two, so $42$ is not a \ealgloss{prime}{liczba pierwsza} number.}

\prq{43}{$1$, $43$}
  {$43$ has exactly two \ealgloss{factors}{dzielniki}, $1$ and itself, so $43$ is a \ealgloss{prime}{liczba pierwsza} number.\\[0.45em]We only need to test $2$, $3$ and $5$ as \ealgloss{factors}{dzielniki}, because $7\times 7=49$ is bigger than $43$.}

\prq{44}{$1$, $2$, $4$, $11$, $22$, $44$}
  {$44$ has six \ealgloss{factors}{dzielniki}, not two, so $44$ is not a \ealgloss{prime}{liczba pierwsza} number.}

\prq{45}{$1$, $3$, $5$, $9$, $15$, $45$}
  {$45$ has six \ealgloss{factors}{dzielniki}, not two, so $45$ is not a \ealgloss{prime}{liczba pierwsza} number.}

\prq{46}{$1$, $2$, $23$, $46$}
  {$46$ has four \ealgloss{factors}{dzielniki}, not two, so $46$ is not a \ealgloss{prime}{liczba pierwsza} number.}

\prq{47}{$1$, $47$}
  {$47$ has exactly two \ealgloss{factors}{dzielniki}, $1$ and itself, so $47$ is a \ealgloss{prime}{liczba pierwsza} number.\\[0.45em]We only need to test $2$, $3$ and $5$ as \ealgloss{factors}{dzielniki}, because $7\times 7=49$ is bigger than $47$.}

\prq{48}{$1$, $2$, $3$, $4$, $6$, $8$, $12$, $16$, $24$, $48$}
  {$48$ has ten \ealgloss{factors}{dzielniki}, not two, so $48$ is not a \ealgloss{prime}{liczba pierwsza} number.}

\prq{49}{$1$, $7$, $49$}
  {$49$ has three \ealgloss{factors}{dzielniki}, not two, so $49$ is not a \ealgloss{prime}{liczba pierwsza} number.}

\prq{50}{$1$, $2$, $5$, $10$, $25$, $50$}
  {$50$ has six \ealgloss{factors}{dzielniki}, not two, so $50$ is not a \ealgloss{prime}{liczba pierwsza} number.}

\prq{51}{$1$, $3$, $17$, $51$}
  {$51$ has four \ealgloss{factors}{dzielniki}, not two, so $51$ is not a \ealgloss{prime}{liczba pierwsza} number.}

\prq{52}{$1$, $2$, $4$, $13$, $26$, $52$}
  {$52$ has six \ealgloss{factors}{dzielniki}, not two, so $52$ is not a \ealgloss{prime}{liczba pierwsza} number.}

\prq{53}{$1$, $53$}
  {$53$ has exactly two \ealgloss{factors}{dzielniki}, $1$ and itself, so $53$ is a \ealgloss{prime}{liczba pierwsza} number.\\[0.45em]We only need to test $2$, $3$, $5$ and $7$ as \ealgloss{factors}{dzielniki}, because $11\times 11=121$ is bigger than $53$.}

\prq{54}{$1$, $2$, $3$, $6$, $9$, $18$, $27$, $54$}
  {$54$ has eight \ealgloss{factors}{dzielniki}, not two, so $54$ is not a \ealgloss{prime}{liczba pierwsza} number.}

\prq{55}{$1$, $5$, $11$, $55$}
  {$55$ has four \ealgloss{factors}{dzielniki}, not two, so $55$ is not a \ealgloss{prime}{liczba pierwsza} number.}

\prq{56}{$1$, $2$, $4$, $7$, $8$, $14$, $28$, $56$}
  {$56$ has eight \ealgloss{factors}{dzielniki}, not two, so $56$ is not a \ealgloss{prime}{liczba pierwsza} number.}

\prq{57}{$1$, $3$, $19$, $57$}
  {$57$ has four \ealgloss{factors}{dzielniki}, not two, so $57$ is not a \ealgloss{prime}{liczba pierwsza} number.}

\prq{58}{$1$, $2$, $29$, $58$}
  {$58$ has four \ealgloss{factors}{dzielniki}, not two, so $58$ is not a \ealgloss{prime}{liczba pierwsza} number.}

\prq{59}{$1$, $59$}
  {$59$ has exactly two \ealgloss{factors}{dzielniki}, $1$ and itself, so $59$ is a \ealgloss{prime}{liczba pierwsza} number.\\[0.45em]We only need to test $2$, $3$, $5$ and $7$ as \ealgloss{factors}{dzielniki}, because $11\times 11=121$ is bigger than $59$.}

\prq{60}{$1$, $2$, $3$, $4$, $5$, $6$, $10$, $12$, $15$, $20$, $30$, $60$}
  {$60$ has twelve \ealgloss{factors}{dzielniki}, not two, so $60$ is not a \ealgloss{prime}{liczba pierwsza} number.}

\prq{61}{$1$, $61$}
  {$61$ has exactly two \ealgloss{factors}{dzielniki}, $1$ and itself, so $61$ is a \ealgloss{prime}{liczba pierwsza} number.\\[0.45em]We only need to test $2$, $3$, $5$ and $7$ as \ealgloss{factors}{dzielniki}, because $11\times 11=121$ is bigger than $61$.}

\prq{62}{$1$, $2$, $31$, $62$}
  {$62$ has four \ealgloss{factors}{dzielniki}, not two, so $62$ is not a \ealgloss{prime}{liczba pierwsza} number.}

\prq{63}{$1$, $3$, $7$, $9$, $21$, $63$}
  {$63$ has six \ealgloss{factors}{dzielniki}, not two, so $63$ is not a \ealgloss{prime}{liczba pierwsza} number.}

\prq{64}{$1$, $2$, $4$, $8$, $16$, $32$, $64$}
  {$64$ has seven \ealgloss{factors}{dzielniki}, not two, so $64$ is not a \ealgloss{prime}{liczba pierwsza} number.}

\prq{65}{$1$, $5$, $13$, $65$}
  {$65$ has four \ealgloss{factors}{dzielniki}, not two, so $65$ is not a \ealgloss{prime}{liczba pierwsza} number.}

\prq{66}{$1$, $2$, $3$, $6$, $11$, $22$, $33$, $66$}
  {$66$ has eight \ealgloss{factors}{dzielniki}, not two, so $66$ is not a \ealgloss{prime}{liczba pierwsza} number.}

\prq{67}{$1$, $67$}
  {$67$ has exactly two \ealgloss{factors}{dzielniki}, $1$ and itself, so $67$ is a \ealgloss{prime}{liczba pierwsza} number.\\[0.45em]We only need to test $2$, $3$, $5$ and $7$ as \ealgloss{factors}{dzielniki}, because $11\times 11=121$ is bigger than $67$.}

\prq{68}{$1$, $2$, $4$, $17$, $34$, $68$}
  {$68$ has six \ealgloss{factors}{dzielniki}, not two, so $68$ is not a \ealgloss{prime}{liczba pierwsza} number.}

\prq{69}{$1$, $3$, $23$, $69$}
  {$69$ has four \ealgloss{factors}{dzielniki}, not two, so $69$ is not a \ealgloss{prime}{liczba pierwsza} number.}

\prq{70}{$1$, $2$, $5$, $7$, $10$, $14$, $35$, $70$}
  {$70$ has eight \ealgloss{factors}{dzielniki}, not two, so $70$ is not a \ealgloss{prime}{liczba pierwsza} number.}

\prq{71}{$1$, $71$}
  {$71$ has exactly two \ealgloss{factors}{dzielniki}, $1$ and itself, so $71$ is a \ealgloss{prime}{liczba pierwsza} number.\\[0.45em]We only need to test $2$, $3$, $5$ and $7$ as \ealgloss{factors}{dzielniki}, because $11\times 11=121$ is bigger than $71$.}

\prq{72}{$1$, $2$, $3$, $4$, $6$, $8$, $9$, $12$, $18$, $24$, $36$, $72$}
  {$72$ has twelve \ealgloss{factors}{dzielniki}, not two, so $72$ is not a \ealgloss{prime}{liczba pierwsza} number.}

\prq{73}{$1$, $73$}
  {$73$ has exactly two \ealgloss{factors}{dzielniki}, $1$ and itself, so $73$ is a \ealgloss{prime}{liczba pierwsza} number.\\[0.45em]We only need to test $2$, $3$, $5$ and $7$ as \ealgloss{factors}{dzielniki}, because $11\times 11=121$ is bigger than $73$.}

\prq{74}{$1$, $2$, $37$, $74$}
  {$74$ has four \ealgloss{factors}{dzielniki}, not two, so $74$ is not a \ealgloss{prime}{liczba pierwsza} number.}

\prq{75}{$1$, $3$, $5$, $15$, $25$, $75$}
  {$75$ has six \ealgloss{factors}{dzielniki}, not two, so $75$ is not a \ealgloss{prime}{liczba pierwsza} number.}

\prq{76}{$1$, $2$, $4$, $19$, $38$, $76$}
  {$76$ has six \ealgloss{factors}{dzielniki}, not two, so $76$ is not a \ealgloss{prime}{liczba pierwsza} number.}

\prq{77}{$1$, $7$, $11$, $77$}
  {$77$ has four \ealgloss{factors}{dzielniki}, not two, so $77$ is not a \ealgloss{prime}{liczba pierwsza} number.}

\prq{78}{$1$, $2$, $3$, $6$, $13$, $26$, $39$, $78$}
  {$78$ has eight \ealgloss{factors}{dzielniki}, not two, so $78$ is not a \ealgloss{prime}{liczba pierwsza} number.}

\prq{79}{$1$, $79$}
  {$79$ has exactly two \ealgloss{factors}{dzielniki}, $1$ and itself, so $79$ is a \ealgloss{prime}{liczba pierwsza} number.\\[0.45em]We only need to test $2$, $3$, $5$ and $7$ as \ealgloss{factors}{dzielniki}, because $11\times 11=121$ is bigger than $79$.}

\prq{80}{$1$, $2$, $4$, $5$, $8$, $10$, $16$, $20$, $40$, $80$}
  {$80$ has ten \ealgloss{factors}{dzielniki}, not two, so $80$ is not a \ealgloss{prime}{liczba pierwsza} number.}

\prq{81}{$1$, $3$, $9$, $27$, $81$}
  {$81$ has five \ealgloss{factors}{dzielniki}, not two, so $81$ is not a \ealgloss{prime}{liczba pierwsza} number.}

\prq{82}{$1$, $2$, $41$, $82$}
  {$82$ has four \ealgloss{factors}{dzielniki}, not two, so $82$ is not a \ealgloss{prime}{liczba pierwsza} number.}

\prq{83}{$1$, $83$}
  {$83$ has exactly two \ealgloss{factors}{dzielniki}, $1$ and itself, so $83$ is a \ealgloss{prime}{liczba pierwsza} number.\\[0.45em]We only need to test $2$, $3$, $5$ and $7$ as \ealgloss{factors}{dzielniki}, because $11\times 11=121$ is bigger than $83$.}

\prq{84}{$1$, $2$, $3$, $4$, $6$, $7$, $12$, $14$, $21$, $28$, $42$, $84$}
  {$84$ has twelve \ealgloss{factors}{dzielniki}, not two, so $84$ is not a \ealgloss{prime}{liczba pierwsza} number.}

\prq{85}{$1$, $5$, $17$, $85$}
  {$85$ has four \ealgloss{factors}{dzielniki}, not two, so $85$ is not a \ealgloss{prime}{liczba pierwsza} number.}

\prq{86}{$1$, $2$, $43$, $86$}
  {$86$ has four \ealgloss{factors}{dzielniki}, not two, so $86$ is not a \ealgloss{prime}{liczba pierwsza} number.}

\prq{87}{$1$, $3$, $29$, $87$}
  {$87$ has four \ealgloss{factors}{dzielniki}, not two, so $87$ is not a \ealgloss{prime}{liczba pierwsza} number.}

\prq{88}{$1$, $2$, $4$, $8$, $11$, $22$, $44$, $88$}
  {$88$ has eight \ealgloss{factors}{dzielniki}, not two, so $88$ is not a \ealgloss{prime}{liczba pierwsza} number.}

\prq{89}{$1$, $89$}
  {$89$ has exactly two \ealgloss{factors}{dzielniki}, $1$ and itself, so $89$ is a \ealgloss{prime}{liczba pierwsza} number.\\[0.45em]We only need to test $2$, $3$, $5$ and $7$ as \ealgloss{factors}{dzielniki}, because $11\times 11=121$ is bigger than $89$.}

\prq{90}{$1$, $2$, $3$, $5$, $6$, $9$, $10$, $15$, $18$, $30$, $45$, $90$}
  {$90$ has twelve \ealgloss{factors}{dzielniki}, not two, so $90$ is not a \ealgloss{prime}{liczba pierwsza} number.}

\prq{91}{$1$, $7$, $13$, $91$}
  {$91$ has four \ealgloss{factors}{dzielniki}, not two, so $91$ is not a \ealgloss{prime}{liczba pierwsza} number.}

\prq{92}{$1$, $2$, $4$, $23$, $46$, $92$}
  {$92$ has six \ealgloss{factors}{dzielniki}, not two, so $92$ is not a \ealgloss{prime}{liczba pierwsza} number.}

\prq{93}{$1$, $3$, $31$, $93$}
  {$93$ has four \ealgloss{factors}{dzielniki}, not two, so $93$ is not a \ealgloss{prime}{liczba pierwsza} number.}

\prq{94}{$1$, $2$, $47$, $94$}
  {$94$ has four \ealgloss{factors}{dzielniki}, not two, so $94$ is not a \ealgloss{prime}{liczba pierwsza} number.}

\prq{95}{$1$, $5$, $19$, $95$}
  {$95$ has four \ealgloss{factors}{dzielniki}, not two, so $95$ is not a \ealgloss{prime}{liczba pierwsza} number.}

\prq{96}{$1$, $2$, $3$, $4$, $6$, $8$, $12$, $16$, $24$, $32$, $48$, $96$}
  {$96$ has twelve \ealgloss{factors}{dzielniki}, not two, so $96$ is not a \ealgloss{prime}{liczba pierwsza} number.}

\prq{97}{$1$, $97$}
  {$97$ has exactly two \ealgloss{factors}{dzielniki}, $1$ and itself, so $97$ is a \ealgloss{prime}{liczba pierwsza} number.\\[0.45em]We only need to test $2$, $3$, $5$ and $7$ as \ealgloss{factors}{dzielniki}, because $11\times 11=121$ is bigger than $97$.}

\prq{98}{$1$, $2$, $7$, $14$, $49$, $98$}
  {$98$ has six \ealgloss{factors}{dzielniki}, not two, so $98$ is not a \ealgloss{prime}{liczba pierwsza} number.}

\prq{99}{$1$, $3$, $9$, $11$, $33$, $99$}
  {$99$ has six \ealgloss{factors}{dzielniki}, not two, so $99$ is not a \ealgloss{prime}{liczba pierwsza} number.}

\prq{100}{$1$, $2$, $4$, $5$, $10$, $20$, $25$, $50$, $100$}
  {$100$ has nine \ealgloss{factors}{dzielniki}, not two, so $100$ is not a \ealgloss{prime}{liczba pierwsza} number.}

%================================================================
\qtopic{\ealgloss{Prime factor decomposition}{Rozkład na czynniki pierwsze}}
%================================================================

\tf[\Large]{$12=2^2\times 3$ is the \ealgloss{prime factor decomposition}{rozkład na czynniki pierwsze} of $12$}
  {\textbf{True.} $2$ and $3$ are both \ealgloss{prime}{liczby pierwsze},
   and $4\times 3=12$.}

\tf[\Large]{$2\times 3\times 4$ is the \ealgloss{prime factor decomposition}{rozkład na czynniki pierwsze} of $24$}
  {\textbf{False.} $4$ is not \ealgloss{prime}{liczba pierwsza}.\\[0.3em]
   $24=2^3\times 3$}

\tf[\Large]{Every whole number greater than $1$ can be written as a \ealgloss{product}{iloczyn}
of \ealgloss{primes}{liczby pierwsze}}
  {\textbf{True.} A \ealgloss{prime}{liczba pierwsza} is its own decomposition, and any other number
   keeps splitting until only \ealgloss{primes}{liczby pierwsze} are left.}

\tf[\Large]{A number can have two different \ealgloss{prime factor decompositions}{rozkłady na czynniki pierwsze}}
  {\textbf{False.} Apart from the order you write them in,
   every number has exactly one decomposition.}

\tf[\Large]{$2^3\times 5=30$}
  {\textbf{False.} $2^3\times 5=8\times 5=40$.\\[0.3em]
   $30=2\times 3\times 5$}

\qq[\LARGE]{Write $12$ as a \ealgloss{product}{iloczyn} of its \ealgloss{prime factors}{czynniki pierwsze}}
  {$12=2\times 2\times 3=2^2\times 3$}

\qq[\LARGE]{Write $18$ as a \ealgloss{product}{iloczyn} of its \ealgloss{prime factors}{czynniki pierwsze}}
  {$18=2\times 3\times 3=2\times 3^2$}

\qq[\LARGE]{Write $60$ as a \ealgloss{product}{iloczyn} of its \ealgloss{prime factors}{czynniki pierwsze}}
  {$60=2\times 2\times 3\times 5=2^2\times 3\times 5$}

\qq[\LARGE]{Write $84$ as a \ealgloss{product}{iloczyn} of its \ealgloss{prime factors}{czynniki pierwsze}}
  {$84=2\times 2\times 3\times 7=2^2\times 3\times 7$}

\qq[\LARGE]{Write $200$ as a \ealgloss{product}{iloczyn} of its \ealgloss{prime factors}{czynniki pierwsze}}
  {$200=2\times 2\times 2\times 5\times 5=2^3\times 5^2$}

\qq[\LARGE]{Write $360$ as a \ealgloss{product}{iloczyn} of its \ealgloss{prime factors}{czynniki pierwsze}}
  {$360=2^3\times 3^2\times 5$}

\qq[\LARGE]{Write $1001$ as a \ealgloss{product}{iloczyn} of its \ealgloss{prime factors}{czynniki pierwsze}}
  {$1001=7\times 143=7\times 11\times 13$}

\qq[\LARGE]{Work out $2^3\times 3\times 5^2$ as an ordinary number}
  {$8\times 3\times 25=600$}

\qq[\Large][\Large]{$n=2^3\times 3\times 5$.\\[0.4em]
Write $2n$ as a \ealgloss{product}{iloczyn} of \ealgloss{prime factors}{czynniki pierwsze}.}
  {Doubling adds one more \ealgloss{factor}{dzielnik} of $2$:\\[0.3em]
   $2n=2^4\times 3\times 5$}

\qq[\Large][\Large]{$720=2^4\times 3^2\times 5$.\\[0.4em]
How many \ealgloss{factors}{dzielniki} does $720$ have?}
  {Each \ealgloss{factor}{dzielnik} uses $0$ to $4$ twos, $0$ to $2$ threes and $0$ or $1$ five:\\[0.3em]
   $5\times 3\times 2=30$ \ealgloss{factors}{dzielniki}}

%================================================================
\qtopic{\ealgloss{HCF}{największy wspólny dzielnik} and \ealgloss{LCM}{najmniejsza wspólna wielokrotność} from lists}
%================================================================

\hcflcm{4}{6}
  {$1$, \hcfmark{2}, $4$}
  {$1$, \hcfmark{2}, $3$, $6$}
  {$4$, $8$, \lcmmark{12}, $16$, $20$,\\[0.45em]
   $24$, $28$, $32$, $36$, $40$}
  {$6$, \lcmmark{12}, $18$, $24$, $30$,\\[0.45em]
   $36$, $42$, $48$, $54$, $60$}

\hcflcm{6}{8}
  {$1$, \hcfmark{2}, $3$, $6$}
  {$1$, \hcfmark{2}, $4$, $8$}
  {$6$, $12$, $18$, \lcmmark{24}, $30$,\\[0.45em]
   $36$, $42$, $48$, $54$, $60$}
  {$8$, $16$, \lcmmark{24}, $32$, $40$,\\[0.45em]
   $48$, $56$, $64$, $72$, $80$}

\hcflcm{6}{9}
  {$1$, $2$, \hcfmark{3}, $6$}
  {$1$, \hcfmark{3}, $9$}
  {$6$, $12$, \lcmmark{18}, $24$, $30$,\\[0.45em]
   $36$, $42$, $48$, $54$, $60$}
  {$9$, \lcmmark{18}, $27$, $36$, $45$,\\[0.45em]
   $54$, $63$, $72$, $81$, $90$}

\hcflcm{8}{12}
  {$1$, $2$, \hcfmark{4}, $8$}
  {$1$, $2$, $3$, \hcfmark{4}, $6$, $12$}
  {$8$, $16$, \lcmmark{24}, $32$, $40$,\\[0.45em]
   $48$, $56$, $64$, $72$, $80$}
  {$12$, \lcmmark{24}, $36$, $48$, $60$,\\[0.45em]
   $72$, $84$, $96$, $108$, $120$}

\hcflcm{10}{15}
  {$1$, $2$, \hcfmark{5}, $10$}
  {$1$, $3$, \hcfmark{5}, $15$}
  {$10$, $20$, \lcmmark{30}, $40$, $50$,\\[0.45em]
   $60$, $70$, $80$, $90$, $100$}
  {$15$, \lcmmark{30}, $45$, $60$, $75$,\\[0.45em]
   $90$, $105$, $120$, $135$, $150$}

\hcflcm{9}{12}
  {$1$, \hcfmark{3}, $9$}
  {$1$, $2$, \hcfmark{3}, $4$, $6$, $12$}
  {$9$, $18$, $27$, \lcmmark{36}, $45$,\\[0.45em]
   $54$, $63$, $72$, $81$, $90$}
  {$12$, $24$, \lcmmark{36}, $48$, $60$,\\[0.45em]
   $72$, $84$, $96$, $108$, $120$}

\hcflcm{12}{18}
  {$1$, $2$, $3$, $4$, \hcfmark{6}, $12$}
  {$1$, $2$, $3$, \hcfmark{6}, $9$, $18$}
  {$12$, $24$, \lcmmark{36}, $48$, $60$,\\[0.45em]
   $72$, $84$, $96$, $108$, $120$}
  {$18$, \lcmmark{36}, $54$, $72$, $90$,\\[0.45em]
   $108$, $126$, $144$, $162$, $180$}

\hcflcm{14}{21}
  {$1$, $2$, \hcfmark{7}, $14$}
  {$1$, $3$, \hcfmark{7}, $21$}
  {$14$, $28$, \lcmmark{42}, $56$, $70$,\\[0.45em]
   $84$, $98$, $112$, $126$, $140$}
  {$21$, \lcmmark{42}, $63$, $84$, $105$,\\[0.45em]
   $126$, $147$, $168$, $189$, $210$}

\hcflcm{15}{20}
  {$1$, $3$, \hcfmark{5}, $15$}
  {$1$, $2$, $4$, \hcfmark{5}, $10$, $20$}
  {$15$, $30$, $45$, \lcmmark{60}, $75$,\\[0.45em]
   $90$, $105$, $120$, $135$, $150$}
  {$20$, $40$, \lcmmark{60}, $80$, $100$,\\[0.45em]
   $120$, $140$, $160$, $180$, $200$}

\hcflcm{8}{9}
  {\hcfmark{1}, $2$, $4$, $8$}
  {\hcfmark{1}, $3$, $9$}
  {$8$, $16$, $24$, $32$, $40$,\\[0.45em]
   $48$, $56$, $64$, \lcmmark{72}, $80$}
  {$9$, $18$, $27$, $36$, $45$,\\[0.45em]
   $54$, $63$, \lcmmark{72}, $81$, $90$}

%================================================================
\qtopic{\ealgloss{HCF}{największy wspólny dzielnik} and \ealgloss{LCM}{najmniejsza wspólna wielokrotność} by \ealgloss{prime factor decomposition}{rozkład na czynniki pierwsze}}
%================================================================

\tf[\Large]{The \ealgloss{HCF}{największy wspólny dzielnik} of two numbers is always smaller than both of them}
  {\textbf{False.} The \ealgloss{HCF}{największy wspólny dzielnik} of $6$ and $12$ is $6$,
   which is one of the numbers itself.}

\tf[\Large]{The \ealgloss{LCM}{najmniejsza wspólna wielokrotność} of two numbers is always a \ealgloss{multiple}{wielokrotność} of both of them}
  {\textbf{True.} That is exactly what a common \ealgloss{multiple}{wielokrotność} is;
   the \ealgloss{LCM}{najmniejsza wspólna wielokrotność} is the smallest one.}

\tf[\Large]{If $a=2^2\times 3$ and $b=2\times 3^2$,\\[0.3em]then the \ealgloss{HCF}{największy wspólny dzielnik} of
$a$ and $b$ is $2\times 3$}
  {\textbf{True.} Take the \emph{lowest} \ealgloss{power}{potęga} of each \ealgloss{prime}{liczba pierwsza} that appears
   in both: $2^1\times 3^1=6$.}

\tf[\Large]{$\mathrm{HCF}\times\mathrm{LCM}$ gives the \ealgloss{product}{iloczyn} of the two numbers}
  {\textbf{True.} Every \ealgloss{prime}{liczba pierwsza} is counted once at its lowest \ealgloss{power}{potęga} and once
   at its highest, which is the same as counting both numbers.}

\tf[\Large]{The \ealgloss{lowest common multiple}{najmniejsza wspólna wielokrotność} of $6$ and $8$ is $48$}
  {\textbf{False.} $6=2\times 3$ and $8=2^3$,
   so the \ealgloss{LCM}{najmniejsza wspólna wielokrotność} is $2^3\times 3=24$.}

\qq[\Large][\Large]{$12=2^2\times 3$ \quad and \quad $18=2\times 3^2$\\[0.4em]
Find the \ealgloss{HCF}{największy wspólny dzielnik} of $12$ and $18$.}
  {Lowest \ealgloss{power}{potęga} of each shared \ealgloss{prime}{liczba pierwsza}:\\[0.3em]
   $\mathrm{HCF}=2\times 3=6$}

\qq[\Large][\Large]{$12=2^2\times 3$ \quad and \quad $18=2\times 3^2$\\[0.4em]
Find the \ealgloss{LCM}{najmniejsza wspólna wielokrotność} of $12$ and $18$.}
  {Highest \ealgloss{power}{potęga} of each \ealgloss{prime}{liczba pierwsza}:\\[0.3em]
   $\mathrm{LCM}=2^2\times 3^2=36$}

\qq[\Large][\Large]{$24=2^3\times 3$ \quad and \quad $60=2^2\times 3\times 5$\\[0.4em]
Find the \ealgloss{HCF}{największy wspólny dzielnik} of $24$ and $60$.}
  {$\mathrm{HCF}=2^2\times 3=12$}

\qq[\Large][\Large]{$24=2^3\times 3$ \quad and \quad $60=2^2\times 3\times 5$\\[0.4em]
Find the \ealgloss{LCM}{najmniejsza wspólna wielokrotność} of $24$ and $60$.}
  {$\mathrm{LCM}=2^3\times 3\times 5=120$}

\qq[\Large][\large]{Find the \ealgloss{HCF}{największy wspólny dzielnik} and the \ealgloss{LCM}{najmniejsza wspólna wielokrotność} of $84$ and $120$}
  {$84=2^2\times 3\times 7$ \qquad $120=2^3\times 3\times 5$\\[0.4em]
   $\mathrm{HCF}=2^2\times 3=12$\\[0.3em]
   $\mathrm{LCM}=2^3\times 3\times 5\times 7=840$}

\qq[\Large][\large]{Find the \ealgloss{HCF}{największy wspólny dzielnik} and the \ealgloss{LCM}{najmniejsza wspólna wielokrotność} of $126$ and $180$}
  {$126=2\times 3^2\times 7$ \qquad $180=2^2\times 3^2\times 5$\\[0.4em]
   $\mathrm{HCF}=2\times 3^2=18$\\[0.3em]
   $\mathrm{LCM}=2^2\times 3^2\times 5\times 7=1260$}

\qq[\Large][\large]{$a=2^3\times 3^2\times 5$ \quad and \quad
$b=2^2\times 3\times 5^2$\\[0.4em]
Find the \ealgloss{HCF}{największy wspólny dzielnik} and the \ealgloss{LCM}{najmniejsza wspólna wielokrotność}, in \ealgloss{index form}{postać potęgowa}.}
  {$\mathrm{HCF}=2^2\times 3\times 5$\\[0.3em]
   $\mathrm{LCM}=2^3\times 3^2\times 5^2$}

\qq[\Large][\large]{Two numbers have \ealgloss{HCF}{największy wspólny dzielnik} $6$ and \ealgloss{LCM}{najmniejsza wspólna wielokrotność} $72$.\\[0.4em]
One of the numbers is $24$. Find the other.}
  {$\mathrm{HCF}\times\mathrm{LCM}=6\times 72=432$, which is the
   \ealgloss{product}{iloczyn} of the two numbers.\\[0.4em]
   $432\div 24=18$}

\qq[\Large][\large]{Find the \ealgloss{HCF}{największy wspólny dzielnik} and the \ealgloss{LCM}{najmniejsza wspólna wielokrotność} of\\[0.4em]
$2^5\times 3^2\times 11$ \quad and \quad $2^3\times 3^4\times 7$\\[0.4em]
Leave your answers in \ealgloss{index form}{postać potęgowa}.}
  {$\mathrm{HCF}=2^3\times 3^2$\\[0.3em]
   $\mathrm{LCM}=2^5\times 3^4\times 7\times 11$}

\qq[\Large][\large]{Two numbers have \ealgloss{HCF}{największy wspólny dzielnik} $12$ and \ealgloss{LCM}{najmniejsza wspólna wielokrotność} $180$.\\[0.4em]
Find the numbers. Is there more than one answer?}
  {Their \ealgloss{product}{iloczyn} is $12\times 180=2160$, and both are \ealgloss{multiples}{wielokrotności} of $12$.\\[0.4em]
   $12$ and $180$ \quad or \quad $36$ and $60$\\[0.3em]
   Both pairs work, so there is more than one answer.}

%================================================================
\qtopic{\ealgloss{Square numbers}{Liczby kwadratowe} from \ealgloss{prime factors}{czynniki pierwsze}}
%================================================================

\tf[\Large]{A number is square exactly when every \ealgloss{power}{potęga} in its \ealgloss{prime factor decomposition}{rozkład na czynniki pierwsze} is even}
  {\textbf{True.} Halving every even \ealgloss{power}{potęga} gives the \ealgloss{square root}{pierwiastek}.}

\tf[\Large]{$2^3\times 3^2$ is a \ealgloss{square number}{liczba kwadratowa}}
  {\textbf{False.} The \ealgloss{power}{potęga} of $2$ is odd,
   so the number cannot be split into two equal halves.}

\tf[\Large]{$2^4\times 5^6$ is a \ealgloss{square number}{liczba kwadratowa}}
  {\textbf{True.} $2^4\times 5^6=\left(2^2\times 5^3\right)^2$}

\tf[\Large]{Every \ealgloss{square number}{liczba kwadratowa} has an even number of \ealgloss{factors}{dzielniki}}
  {\textbf{False.} \ealgloss{Square numbers}{Liczby kwadratowe} have an \emph{odd} number of \ealgloss{factors}{dzielniki}:\\[0.3em]
   $16$ has $1,\ 2,\ 4,\ 8,\ 16$.}

\tf[\Large]{If $n$ is a \ealgloss{square number}{liczba kwadratowa}, then $4n$ is a \ealgloss{square number}{liczba kwadratowa} too}
  {\textbf{True.} Multiplying by $4=2^2$ keeps every \ealgloss{power}{potęga} even.}

\qq[\Large][\Large]{Use \ealgloss{prime factors}{czynniki pierwsze} to show that $36$ is a \ealgloss{square number}{liczba kwadratowa}}
  {$36=2^2\times 3^2$\\[0.3em]
   Both \ealgloss{powers}{potęgi} are even, and $36=(2\times 3)^2=6^2$.}

\qq[\Large][\Large]{Is $2^2\times 3^2\times 5^2$ a \ealgloss{square number}{liczba kwadratowa}?\\[0.4em]
If so, what is its \ealgloss{square root}{pierwiastek}?}
  {Yes --- every \ealgloss{power}{potęga} is even.\\[0.3em]
   $\sqrt{2^2\times 3^2\times 5^2}=2\times 3\times 5=30$}

\qq[\Large][\Large]{Is $2^3\times 3^2$ a \ealgloss{square number}{liczba kwadratowa}?}
  {No. Halving the \ealgloss{powers}{potęgi} would need $2^{1.5}$,\\[0.3em]
   so the \ealgloss{power}{potęga} of $2$ must be even and it is not.}

\qq[\Large][\Large]{Work out $\sqrt{2^6\times 3^4}$}
  {Halve each \ealgloss{power}{potęga}:\\[0.3em]
   $2^3\times 3^2=8\times 9=72$}

\qq[\Large][\Large]{Use \ealgloss{prime factors}{czynniki pierwsze} to decide whether $400$ is a \ealgloss{square number}{liczba kwadratowa}}
  {$400=2^4\times 5^2$\\[0.3em]
   Both \ealgloss{powers}{potęgi} are even, so $400=\left(2^2\times 5\right)^2=20^2$.}

\qq[\Large][\Large]{Use \ealgloss{prime factors}{czynniki pierwsze} to decide whether $150$ is a \ealgloss{square number}{liczba kwadratowa}}
  {$150=2\times 3\times 5^2$\\[0.3em]
   The \ealgloss{powers}{potęgi} of $2$ and $3$ are odd, so $150$ is not square.}

\qq[\Large][\Large]{Work out $\sqrt{2^8\times 3^2\times 5^4}$}
  {Halve each \ealgloss{power}{potęga}:\\[0.3em]
   $2^4\times 3\times 5^2=16\times 3\times 25=1200$}

\qq[\Large][\large]{$n=2^3\times 3\times 5^2$\\[0.4em]
Find the smallest whole number $k$ so that $nk$ is a \ealgloss{square number}{liczba kwadratowa}.}
  {The \ealgloss{powers}{potęgi} of $2$ and $3$ are odd, so each needs one more:\\[0.3em]
   $k=2\times 3=6$, giving $nk=2^4\times 3^2\times 5^2$.}

\qq[\Large][\large]{$504=2^3\times 3^2\times 7$\\[0.4em]
Find the smallest number $504$ can be multiplied by to give a \ealgloss{square number}{liczba kwadratowa}.}
  {The \ealgloss{powers}{potęgi} of $2$ and $7$ are odd:\\[0.3em]
   multiply by $2\times 7=14$.\\[0.3em]
   $504\times 14=7056=84^2$}

\qq[\Large][\large]{$1080=2^3\times 3^3\times 5$\\[0.4em]
Find the smallest number $1080$ can be divided by to give a \ealgloss{square number}{liczba kwadratowa}.}
  {Remove one $2$, one $3$ and the $5$:\\[0.3em]
   divide by $2\times 3\times 5=30$.\\[0.3em]
   $1080\div 30=36=6^2$}

%================================================================
\qtopic{Problem solving with \ealgloss{HCF}{największy wspólny dzielnik} and \ealgloss{LCM}{najmniejsza wspólna wielokrotność}}
%================================================================

\tf[\Large]{To find when two repeating events next happen together,
you use the \ealgloss{LCM}{najmniejsza wspólna wielokrotność}}
  {\textbf{True.} You are looking for the first time that is a \ealgloss{multiple}{wielokrotność}
   of both gaps.}

\tf[\Large]{To split two collections into the largest identical bags,
you use the \ealgloss{LCM}{najmniejsza wspólna wielokrotność}}
  {\textbf{False.} Use the \ealgloss{HCF}{największy wspólny dzielnik} --- the bag size has to be a \ealgloss{factor}{dzielnik}
   of both amounts.}

\tf[\Large]{Cutting two ropes into equal pieces, as long as possible with
none left over, uses the \ealgloss{HCF}{największy wspólny dzielnik}}
  {\textbf{True.} The piece length must be a \ealgloss{factor}{dzielnik} of both rope lengths.}

\tf[\Large]{The \ealgloss{LCM}{najmniejsza wspólna wielokrotność} of two numbers is never smaller than either of them}
  {\textbf{True.} It is a \ealgloss{multiple}{wielokrotność} of each number,
   so it is at least as big as both.}

\tf[\Large]{Two events repeat every $4$ minutes and every $6$ minutes,\\[0.3em]
so they coincide every $24$ minutes}
  {\textbf{False.} They coincide every $\mathrm{LCM}(4,6)=12$ minutes.}

\qq[\Large][\Large]{Two bells ring together now.\\[0.4em]
One rings every $4$ minutes, the other every $6$ minutes.\\[0.4em]
When do they next ring together?}
  {Common \ealgloss{multiple}{wielokrotność}, so use the \ealgloss{LCM}{najmniejsza wspólna wielokrotność}:\\[0.3em]
   $\mathrm{LCM}(4,6)=12$ --- in $12$ minutes.}

\qq[\Large][\Large]{Two ribbons, $18$\,cm and $24$\,cm long, are cut into
equal pieces, as long as possible with none left over.\\[0.4em]
How long is each piece?}
  {The length must be a \ealgloss{factor}{dzielnik} of both, so use the \ealgloss{HCF}{największy wspólny dzielnik}:\\[0.3em]
   $\mathrm{HCF}(18,24)=6$ --- each piece is $6$\,cm.}

\qq[\Large][\large]{$36$ pens and $48$ pencils are shared into identical packs
with nothing left over.\\[0.4em]
What is the greatest number of packs?}
  {$36=2^2\times 3^2$ \qquad $48=2^4\times 3$\\[0.4em]
   $\mathrm{HCF}=2^2\times 3=12$ packs,\\[0.3em]
   each with $3$ pens and $4$ pencils.}

\qq[\Large][\Large]{Two lighthouses flash together at $8$pm.\\[0.4em]
One flashes every $15$ seconds, the other every $18$ seconds.\\[0.4em]
When do they next flash together?}
  {$\mathrm{LCM}(15,18)=90$ seconds\\[0.3em]
   At $8${:}$01${:}$30$pm.}

\qq[\Large][\Large]{Rectangular tiles measure $12$\,cm by $18$\,cm.\\[0.4em]
What is the smallest square they can tile exactly?}
  {The side must be a \ealgloss{multiple}{wielokrotność} of both, so use the \ealgloss{LCM}{najmniejsza wspólna wielokrotność}:\\[0.3em]
   $\mathrm{LCM}(12,18)=36$\\[0.3em]
   A $36$\,cm by $36$\,cm square.}

\qq[\Large][\Large]{A bus leaves every $24$ minutes and a tram every
$36$ minutes.\\[0.4em]
Both leave at $9${:}$00$. When do they next leave together?}
  {$\mathrm{LCM}(24,36)=72$ minutes\\[0.3em]
   At $10${:}$12$.}

\qq[\Large][\large]{Two cogs, with $18$ and $24$ teeth, start meshed at a
marked tooth.\\[0.4em]
How many turns does each make before the marks meet again?}
  {$\mathrm{LCM}(18,24)=72$ teeth must pass.\\[0.4em]
   $72\div 18=4$ turns of the small cog,\\[0.3em]
   $72\div 24=3$ turns of the large cog.}

\qq[\Large][\large]{A floor measures $360$\,cm by $504$\,cm.\\[0.4em]
It is tiled with identical square tiles, as large as possible.\\[0.4em]
Find the tile size and the number of tiles.}
  {$360=2^3\times 3^2\times 5$ \qquad $504=2^3\times 3^2\times 7$\\[0.4em]
   $\mathrm{HCF}=2^3\times 3^2=72$, so the tiles are $72$\,cm square.\\[0.3em]
   $5\times 7=35$ tiles.}

\qq[\Large][\large]{Three runners take $60$\,s, $72$\,s and $90$\,s per lap.\\[0.4em]
They start together. When are they next all at the start line?}
  {$60=2^2\times 3\times 5$ \quad $72=2^3\times 3^2$ \quad
   $90=2\times 3^2\times 5$\\[0.4em]
   $\mathrm{LCM}=2^3\times 3^2\times 5=360$\,s\\[0.3em]
   After $6$ minutes.}

\qq[\Large][\large]{Hot dogs come in packs of $10$, buns in packs of $8$ and
sauce sachets in packs of $12$.\\[0.4em]
What is the smallest number of complete hot dogs Sam can make with nothing
left over, and how many packs of each is that?}
  {$\mathrm{LCM}(10,8,12)=2^3\times 3\times 5=120$\\[0.4em]
   $120$ hot dogs: $12$ packs of hot dogs,\\[0.3em]
   $15$ packs of buns and $10$ packs of sachets.}

\end{document}